\chapter{Angular Momentum, Central Forces, and the Oscillator Ladder}

These notes follow Leonard Susskind's \emph{Advanced Quantum Mechanics} Lecture~3, curated by LazyingArt LLC. The lecture has a clear internal rhythm. Angular momentum is first reintroduced as a statement about the angular dependence of a wavefunction, then used to reduce the central-force Schr\"odinger equation to an effective radial problem, and finally echoed in the ladder-operator treatment of the harmonic oscillator. Several factors of \(\hbar\) are suppressed on the blackboard; they are restored here in standard form.

\section{Angular Momentum as Angular Dependence}

Susskind begins not with operator algebra in the abstract, but with a particle moving in a central force field. Classically, the angular momentum vector is conserved, so the motion stays in a fixed plane. Quantum mechanically, the particle is described by a wavefunction, and the essential point is that angular momentum concerns the \emph{angular} dependence of that wavefunction.

The cleanest prototype is two-dimensional motion. In that case the only rotation is about the axis perpendicular to the plane, and the corresponding generator is
\begin{equation}
L_z=-i\hbar\,\frac{\partial}{\partial \theta}.
\end{equation}
To find angular-momentum eigenstates one solves
\begin{equation}
L_z\psi(r,\theta)=\hbar \ell\,\psi(r,\theta),
\end{equation}
or explicitly
\begin{equation}
-i\hbar\,\frac{\partial}{\partial \theta}\psi(r,\theta)=\hbar \ell\,\psi(r,\theta).
\end{equation}
Since the derivative is with respect to \(\theta\) alone, the radial coordinate \(r\) is merely a spectator. The solution is therefore
\begin{equation}
\psi(r,\theta)=\chi(r)e^{i\ell\theta},
\end{equation}
with \(\ell\) an integer so that the wavefunction is single-valued when \(\theta\) advances by \(2\pi\).

This is the first important structural lesson of the lecture: angular momentum is encoded in the angular part of the wavefunction, while the radial factor is irrelevant to the angular-momentum eigenvalue.

In three dimensions the same logic survives. One writes the wavefunction schematically as
\begin{equation}
\Psi(r,\theta,\phi)=R(r)\,Y_{\ell m}(\theta,\phi),
\end{equation}
where the spherical harmonics \(Y_{\ell m}\) carry the rotational quantum numbers. Susskind's point is not that one must know their explicit formulas immediately, but that the angular problem and the radial problem can be separated conceptually from the start.

\section{Ladder Operators and the Structure of Multiplets}

Having reminded the audience what angular momentum is about, the lecture returns to the algebra from the previous meeting. One starts with the familiar commutators
\begin{equation}
[L_x,L_y]=i\hbar L_z,\qquad
[L_y,L_z]=i\hbar L_x,\qquad
[L_z,L_x]=i\hbar L_y.
\end{equation}
The Dirac-inspired trick is to combine \(L_x\) and \(L_y\) into
\begin{equation}
L_\pm=L_x\pm iL_y.
\end{equation}
Working in an \(L_z\)-eigenbasis,
\begin{equation}
L_z|\ell,m\rangle=\hbar m\,|\ell,m\rangle,
\end{equation}
one finds that \(L_+\) and \(L_-\) raise and lower the \(m\)-value by one unit:
\begin{equation}
L_\pm|\ell,m\rangle \propto |\ell,m\pm 1\rangle.
\end{equation}
Thus the spectrum of \(m\) is arranged into ladders.

The next step in the lecture is qualitative but decisive. If repeated application of \(L_+\) never ended, one would obtain an infinite ladder. A finite multiplet appears only when the ladder terminates above and below:
\begin{equation}
L_+|\ell,\ell\rangle=0,\qquad
L_-|\ell,-\ell\rangle=0.
\end{equation}
Rotational symmetry then forces the ladder to be symmetric about \(m=0\), so that
\begin{equation}
m=-\ell,-\ell+1,\dots,\ell-1,\ell.
\end{equation}
The number of states in such a multiplet is therefore
\begin{equation}
2\ell+1.
\end{equation}

\begin{remark}
The ladder argument by itself permits both integer and half-integer multiplets. Later in the lecture, when the states are interpreted as orbital wavefunctions \(Y_{\ell m}(\theta,\phi)\) on the sphere, \(\ell\) is restricted to nonnegative integers.
\end{remark}

\begin{figure}
\centering
\begin{tikzpicture}[x=1cm,y=0.9cm]
\draw[->] (0,-3.5) -- (0,3.5) node[above] {$m$};

\fill (0,3) circle (2pt);
\fill (0,2) circle (2pt);
\fill (0,0) circle (2pt);
\fill (0,-2) circle (2pt);
\fill (0,-3) circle (2pt);

\node[right] at (0,3) {$\ell$};
\node[right] at (0,2) {$\ell-1$};
\node[right] at (0,0) {$0$};
\node[right] at (0,-2) {$-\ell+1$};
\node[right] at (0,-3) {$-\ell$};

\node at (0,1) {$\vdots$};
\node at (0,-1) {$\vdots$};

\draw[->] (-1.2,-2.4) -- (-1.2,2.4) node[midway,left] {$L_+$};
\draw[->] (1.2,2.4) -- (1.2,-2.4) node[midway,right] {$L_-$};
\end{tikzpicture}
\caption{A finite angular-momentum multiplet. The operators \(L_+\) and \(L_-\) move upward and downward in integer steps of \(m\), and the ladder terminates at \(m=\pm \ell\).}
\end{figure}

\section{The Common Value of \(L^2\)}

The algebraic climax of the first half of the lecture is the determination of the common \(L^2\)-eigenvalue for the entire multiplet. One begins from
\begin{equation}
L^2=L_x^2+L_y^2+L_z^2
\end{equation}
and rewrites the transverse part using the ladder operators. The ordered product
\begin{align}
L_-L_+&=(L_x-iL_y)(L_x+iL_y) \\
&=L_x^2+L_y^2+i[L_x,L_y] \\
&=L_x^2+L_y^2-\hbar L_z.
\end{align}
Hence
\begin{equation}
L^2=L_z^2+\hbar L_z+L_-L_+.
\end{equation}
Equivalently,
\begin{equation}
L^2=L_z^2-\hbar L_z+L_+L_-.
\end{equation}
Both forms are useful; the first is adapted to the top state, the second to the bottom state.

\paragraph{Worked derivation of the \(\ell(\ell+1)\) law.}
Apply the first identity to the highest state of the multiplet:
\begin{align}
L^2|\ell,\ell\rangle
&=\left(L_z^2+\hbar L_z+L_-L_+\right)|\ell,\ell\rangle \\
&=\left(\hbar^2\ell^2+\hbar^2\ell\right)|\ell,\ell\rangle
+L_-\bigl(L_+|\ell,\ell\rangle\bigr).
\end{align}
But \(L_+|\ell,\ell\rangle=0\) by definition of the top state, so
\begin{equation}
L^2|\ell,\ell\rangle=\hbar^2\ell(\ell+1)|\ell,\ell\rangle.
\end{equation}
This is the blackboard result summarized in the lecture by the shorthand \(L^2=\ell(\ell+1)\), with the factors of \(\hbar^2\) understood.

The lecture then immediately adds the second important point: every state in the multiplet has the same value of \(L^2\). The reason is that \(L^2\) commutes with every component of angular momentum, hence with \(L_\pm\). Once the top state has eigenvalue \(\hbar^2\ell(\ell+1)\), repeated lowering preserves it throughout the ladder.

Thus the orbital angular states are labeled by two quantum numbers:
\begin{equation}
L^2Y_{\ell m}=\hbar^2\ell(\ell+1)Y_{\ell m},\qquad
L_zY_{\ell m}=\hbar m\,Y_{\ell m},
\end{equation}
with \(m=-\ell,\dots,\ell\). In a central force problem these states are degenerate within a fixed \(\ell\)-multiplet because rotational invariance implies that the Hamiltonian commutes with angular momentum.

\section{Central Forces and the Effective Radial Hamiltonian}

Only after the angular-momentum structure is secure does the lecture turn back to the central-force Hamiltonian. One starts from
\begin{equation}
H=\frac{p^2}{2m}+V(r),
\end{equation}
where \(V(r)\) depends only on the distance from the origin. Classically, rotational symmetry implies conservation of
\begin{equation}
\mathbf{L}=\mathbf{r}\times\mathbf{p}.
\end{equation}
Because \(\mathbf{L}\) is conserved as a vector, the orbital plane is fixed. One may therefore decompose the momentum in that plane into a radial part \(p_r\) and an angular part \(p_\theta\):
\begin{equation}
p^2=p_r^2+p_\theta^2.
\end{equation}
The Hamiltonian becomes
\begin{equation}
H=\frac{p_r^2+p_\theta^2}{2m}+V(r).
\end{equation}
Now use the classical relation
\begin{equation}
L=r\,p_\theta,
\end{equation}
or
\begin{equation}
L^2=r^2p_\theta^2,\qquad p_\theta^2=\frac{L^2}{r^2}.
\end{equation}
Substituting into the Hamiltonian gives
\begin{equation}
H=\frac{p_r^2}{2m}+\frac{L^2}{2mr^2}+V(r).
\end{equation}
This is the effective one-dimensional radial Hamiltonian emphasized on the blackboard.

The new term
\begin{equation}
\frac{L^2}{2mr^2}
\end{equation}
is the centrifugal barrier. It is repulsive, grows rapidly as \(r\to 0\), and prevents a particle with nonzero angular momentum from getting arbitrarily close to the origin. The only exceptional case is \(\ell=0\), where the barrier disappears.

\paragraph{Physical interpretation.}
For a Coulomb potential \(V(r)\sim -1/r\), the centrifugal term behaves as \(+1/r^2\). Near the origin the barrier dominates, while far away the Coulomb tail dominates. The resulting effective potential has a minimum and supports circular motion at the bottom, together with radial oscillations about that minimum. This is the classical bridge to the quantum radial problem.

\section{The Radial Schr\"odinger Equation and Bound States}

The quantum version of the same story begins from the observation that once the angular momentum is fixed, the angular dependence has already been solved. The angular part of the wavefunction is one of the spherical harmonics \(Y_{\ell m}\), so the remaining unknown is a radial function. In the lecture this radial unknown is written simply as \(\psi(r)\).

Quantization replaces the radial momentum by
\begin{equation}
p_r=-i\hbar\,\frac{d}{dr},
\end{equation}
so that
\begin{equation}
p_r^2\longrightarrow -\hbar^2\,\frac{d^2}{dr^2}.
\end{equation}
At the same time the classical \(L^2\) is replaced, on angular-momentum eigenstates, by \(\hbar^2\ell(\ell+1)\). The radial equation is therefore
\begin{equation}
-\frac{\hbar^2}{2m}\frac{d^2}{dr^2}\psi(r)
+\frac{\hbar^2\ell(\ell+1)}{2mr^2}\psi(r)
+V(r)\psi(r)
=
E\psi(r).
\end{equation}

\begin{remark}
On the blackboard Susskind writes the radial equation in the one-variable form above. To keep the lecture's notation while remaining mathematically clean, one may interpret \(\psi(r)\) here as the reduced radial wavefunction \(u_\ell(r)=rR_\ell(r)\). In terms of the unreduced radial function \(R_\ell(r)\), the separated equation contains an additional first-derivative term.
\end{remark}

This is the main technical payoff of the angular-momentum discussion. The original three-dimensional Schr\"odinger equation,
\begin{equation}
-\frac{\hbar^2}{2m}\left(\frac{\partial^2}{\partial x^2}
+\frac{\partial^2}{\partial y^2}
+\frac{\partial^2}{\partial z^2}\right)\Psi
+V(r)\Psi
=
E\Psi,
\end{equation}
has been reduced to an ordinary differential equation in \(r\).

The lecture does not solve this equation explicitly. Instead it organizes the qualitative spectrum. For fixed \(\ell\), one has a one-dimensional bound-state problem, and the states are ordered by the number of nodes in the radial wavefunction:
\begin{enumerate}
\item the lowest bound state has no nodes,
\item the first excited state has one node,
\item the next has two nodes,
\item and so on.
\end{enumerate}
Susskind's physical explanation is that more nodes mean more rapid oscillation, hence larger momentum and larger kinetic energy.

Each radial level at fixed \(\ell\) carries degeneracy
\begin{equation}
2\ell+1,
\end{equation}
coming from the allowed values of \(m\). For a generic attractive central potential, the \(\ell=0\) tower sits lowest, the \(\ell=1\) tower sits somewhat higher, the \(\ell=2\) tower higher still, and each tower contains its own sequence of node-counted bound states.

There is also a useful special remark in the lecture about \(\ell=0\). Classically, if the centrifugal barrier vanishes, the particle would simply fall to the center. Quantum mechanically, however, one cannot have both sharply localized position at the origin and zero momentum. The ground state is therefore spread over a finite region rather than collapsing to a point.

\section{Generic Atomic Spectra and the Coulomb Accident}

Susskind next arranges the spectrum by plotting angular momentum horizontally and energy vertically. For each value of \(\ell\), the radial equation is different because the centrifugal term changes. One therefore obtains a separate radial tower for each \(\ell\), and the levels are not equally spaced. The only universal features are the node ordering and the multiplicity \(2\ell+1\).

The Coulomb potential \(V(r)\propto -1/r\) is special. The lecture emphasizes that the specialness is real but also fragile. It is exact for the idealized nonrelativistic \(1/r\) problem, but not for real atoms, where finite nuclear size, relativistic corrections, and spin all spoil it.

The special feature is an unexpected alignment of levels with different \(\ell\). In the Coulomb problem,
\begin{itemize}
\item the lowest \(\ell=1\) state lines up with the first excited \(\ell=0\) state,
\item the lowest \(\ell=2\) state lines up with the second excited \(\ell=0\) state,
\item the lowest \(\ell=3\) state lines up with the third excited \(\ell=0\) state,
\item and similarly for the higher towers.
\end{itemize}
The resulting shell degeneracies are
\begin{equation}
1,\quad 4,\quad 9,\quad 16,\quad \dots
\end{equation}
In other words, the \(n\)-th Coulomb shell has orbital degeneracy
\begin{equation}
\sum_{\ell=0}^{n-1}(2\ell+1)=n^2.
\end{equation}
This is precisely the sequence Susskind builds on the board by counting states level by level.

He does not pause to prove the hidden symmetry behind this degeneracy. The lecture treats it as a remarkable and highly special fact about the Coulomb potential, and that is the right emphasis here as well. For hydrogen this is an excellent first approximation, but not the whole story. Once spin is included, each orbital state doubles, and the gross degeneracies become
\begin{equation}
2,\quad 8,\quad 18,\quad 32,\quad \dots
\end{equation}
in agreement with the familiar atomic shell counting.

A final qualitative remark from the discussion is worth preserving: bound-state wavefunctions do not have compact support. They do not vanish exactly beyond some finite radius, but for bound states they decrease rapidly, typically exponentially, so that the tail is practically negligible far from the atom.

\section{The Harmonic Oscillator: A New Ladder}

The lecture then changes subject without changing method. The harmonic oscillator is introduced as the second great example in which Dirac-style algebra works beautifully.

Set the mass equal to one and write the spring constant as \(k=\omega^2\). The classical Hamiltonian is
\begin{equation}
H=\frac{p^2}{2}+\frac{\omega^2x^2}{2}.
\end{equation}
Hamilton's equations give
\begin{equation}
\dot x=\frac{\partial H}{\partial p}=p,\qquad
\dot p=-\frac{\partial H}{\partial x}=-\omega^2x,
\end{equation}
hence
\begin{equation}
\ddot x=-\omega^2x.
\end{equation}
So \(\omega\) is indeed the oscillation frequency.

Now comes the factorization trick. One tries to write the Hamiltonian as a product of first-order operators:
\begin{equation}
\frac12(p+i\omega x)(p-i\omega x).
\end{equation}
Quantum mechanically this is not exactly \(H\), because \(x\) and \(p\) do not commute. Using
\begin{equation}
[x,p]=i\hbar,
\end{equation}
one finds
\begin{align}
\frac12(p+i\omega x)(p-i\omega x)
&=\frac12\bigl(p^2+\omega^2x^2+i\omega[x,p]\bigr) \\
&=\frac12\bigl(p^2+\omega^2x^2-\hbar\omega\bigr).
\end{align}
Therefore
\begin{equation}
H=\frac12(p+i\omega x)(p-i\omega x)+\frac12\hbar\omega.
\end{equation}
The extra constant is the zero-point energy.

Introduce the normalized operators
\begin{equation}
A_+=\frac{p+i\omega x}{\sqrt{2\hbar\omega}},\qquad
A_-=\frac{p-i\omega x}{\sqrt{2\hbar\omega}}.
\end{equation}
They are Hermitian conjugates of each other, and the Hamiltonian becomes
\begin{equation}
H=\hbar\omega\left(A_+A_-+\frac12\right).
\end{equation}
It is convenient to define the number operator
\begin{equation}
N=A_+A_-.
\end{equation}
A short commutator calculation yields
\begin{equation}
[A_-,A_+]=1,
\end{equation}
and therefore
\begin{equation}
[N,A_+]=A_+,\qquad [N,A_-]=-A_-.
\end{equation}
If
\begin{equation}
N|n\rangle=n|n\rangle,
\end{equation}
then
\begin{align}
N(A_+|n\rangle)
&=(A_+N+[N,A_+])|n\rangle \\
&=(n+1)(A_+|n\rangle),
\end{align}
so \(A_+\) raises the eigenvalue by one. Similarly,
\begin{equation}
N(A_-|n\rangle)=(n-1)(A_-|n\rangle).
\end{equation}

This reproduces the angular-momentum logic almost exactly, but now the ladder is a ladder of energies. Because
\begin{equation}
N=A_+A_-
\end{equation}
is positive semidefinite, its eigenvalues cannot decrease indefinitely. There must be a bottom state \(|0\rangle\) satisfying
\begin{equation}
A_-|0\rangle=0.
\end{equation}
With the conventional normalization,
\begin{equation}
A_+|n\rangle=\sqrt{n+1}\,|n+1\rangle,\qquad
A_-|n\rangle=\sqrt{n}\,|n-1\rangle,
\end{equation}
and the spectrum is
\begin{equation}
E_n=\hbar\omega\left(n+\frac12\right),\qquad n=0,1,2,\dots
\end{equation}

The lecture ends just as the normalization factors \(\sqrt{n+1}\) and \(\sqrt{n}\) are being restored. The main structural point, however, is already complete: once again a commutator algebra produces a ladder, a lowest state, and a discrete spectrum.

\section{Summary}

This lecture develops one method across two seemingly different problems. Angular momentum is first understood as the angular part of a wavefunction, organized into finite multiplets by the ladder operators \(L_\pm\), with the characteristic quantum result \(L^2=\hbar^2\ell(\ell+1)\). That algebra is then used to reduce the central-force Schr\"odinger equation to a radial ordinary differential equation with an effective centrifugal barrier. The resulting bound states are ordered by nodes, and the Coulomb potential stands out by its special \(1,4,9,16,\dots\) degeneracy pattern. In the second act, the harmonic oscillator repeats the same algebraic story: factorization, commutator, ladder, ground state, and evenly spaced spectrum.